\subsection{平稳序列}

定义5. 一个序列 \((x_{n})_{n\in \mathbb{N}}\) （集合 \(\mathbf{E}\) 的元素的序列）被称为平稳的，如果存在一个整数 \(m\) 使得 \(x_{n} = x_{m}\) 对于所有整数 \(n\geqslant m\) .

命题6. 设E是一个有序集。则以下陈述等价：

\begin{enumerate}
\def\labelenumi{(\alph{enumi})}
\item
  \(\mathbf{E}\) 的每个非空子集都有一个极大元。
\item
  每个递增序列 \((x_{n})\) （由 E 中的元素组成）都是平稳的。
\end{enumerate}

我们首先证明 (a) 蕴含 (b)。设 X 为序列 \((x_{n})\) 的元素所构成的集合，并且让 \(x_{m}\) 为其一个极大元。若 \(n \geqslant m\) ，根据假设我们有 \(x_{n} \geqslant x_{m}\) ，从而 \(x_{n} = x_{m}\) ，这是根据 \(x_{m}\) 的极大性。反之，假设存在E的一个非空子集A，它没有极大元素。对于每个 \(x \in A\) ，让 \(T_{x}\) 为 \(y \in A\) 中所有使得 y \textgreater{} x 的元素所成的集合。根据假设， \(T_{x} \neq \emptyset\) 对于所有 \(x \in A\) ；因此存在一个从A到A的映射f，使得 \(f(x) > x\) 对于所有 \(x \in A\) （第二章，第5节，第4号，命题6）。如果 \(a \in A\) ，则序列 \((x_{n})_{n \in \mathbb{N}}\) 由以下条件归纳定义： \(x_{0} = a\) , \(x_{n+1} = f(x_{n})\) ，显然递增且非平稳。

推论1. 一个全序集 E 是良序的，当且仅当 E 中元素的每个递减序列都是平稳的。

因为说 E 是良序的等价于说 E 的每个非空子集都有一个极小元素（ \(\S\) 1，第10号，命题10），因此该断言由命题6得出。

推论2. 有限有序集元素的每个递增序列都是平稳的。

因为每个有限有序集都有一个极大元（§4，第4号，命题3，推论2）。

¶ 满足命题6的等价条件的有序集E有时被称为诺特集。

命题7（“诺特归纳原理”）。设E是一个诺特集，并设F是E的一个具有如下性质的子集：如果 \(a \in E\) 使得关系 \(x > a\) 蕴含 \(x \in F\) ，那么 \(a \in F\) 。在这些条件下， \(F = E\) 。

确实，假设 \(\mathbf{E} \neq \mathbf{F}\) ；然后 \(\mathbf{E} - \mathbf{F}\) 具有一个极大元素 \(b\) 。根据定义，我们有 \(x \in \mathbf{F}\) 对于所有 \(x > b\) ；但这意味着 \(b \in \mathbf{F}\) ，这是荒谬的。

